\begin{explanationbox}

	\textbf{\textcolor{CExplanation}{一句话直觉}}

	离心率就是圆锥曲线的“形状指标”：它用“到焦点的距离”和“到准线的距离”的固定比例，统一描述椭圆、抛物线和双曲线。

	\medskip

	\textbf{\textcolor{CExplanation}{核心拆解}}

	\begin{description}[style=nextline, leftmargin=2.8em, labelsep=0.5em, itemsep=0.45em]

		\item[\textbf{要点一（几何比例）：}]
		      对于圆锥曲线上任意一点$P$，它到对应焦点$F$和到对应准线$l$的距离之比是一个常数，这个常数就是离心率$e$。

		      \[
			      e=\frac{|PF|}{d(P,l)}.
		      \]

		\item[\textbf{要点二（代数连接）：}]
		      在椭圆和双曲线中，$e=\frac{c}{a}$，其中$c$表示半焦距，$a$表示半长轴或半实轴长。

		\item[\textbf{要点三（形状指标）：}]
		      $e$的大小决定曲线类型：$0<e<1$为椭圆，$e=1$为抛物线，$e>1$为双曲线。

	\end{description}

	\medskip

	\textbf{\textcolor{CExplanation}{几何本质（曲线形状）}}

	离心率反映曲线偏离“圆形”的程度。对椭圆而言，$e$越接近$0$，椭圆越接近圆；$e$越接近$1$，椭圆越扁。对双曲线而言，$e$越大，标准双曲线的渐近线张角越大，开口趋势越明显。

	\medskip

	\textbf{\textcolor{CExplanation}{代数意义（参数联系）}}

	在椭圆和双曲线的标准方程中，$a$、$b$、$c$通过$e=\frac{c}{a}$以及$c^{2}=a^{2}\pm b^{2}$相互关联。椭圆用$c^{2}=a^{2}-b^{2}$，双曲线用$c^{2}=a^{2}+b^{2}$。

	\medskip

	\textbf{\textcolor{CExplanation}{考点价值（高频应用）}}

	\begin{itemize}[leftmargin=*, itemsep=0.5em, topsep=0.4em]

		\item \textbf{考法一（参数求解）：} 给出椭圆或双曲线标准方程，直接求$e$；或已知$e$和某些参数，反求方程。

		\item \textbf{考法二（定义迁移）：} 题目中出现焦点、准线、距离比例时，优先考虑用定义$e=\frac{|PF|}{d(P,l)}$列等式。

	\end{itemize}

	\medskip

	\textbf{\textcolor{CExplanation}{顿悟点}}

	若焦点在$x$轴上，椭圆或双曲线的准线常写成$x=\pm\frac{a^{2}}{c}$，也可以写成$x=\pm\frac{a}{e}$。如果焦点在$y$轴上，则准线应写成$y=\pm\frac{a^{2}}{c}$或$y=\pm\frac{a}{e}$。

	\medskip

	\textbf{\textcolor{CExplanation}{使用场景}}

	\begin{itemize}[leftmargin=*, itemsep=0.5em, topsep=0.4em]

		\item \textbf{场景一（计算问题）：} 已知椭圆或双曲线标准方程，需要求离心率时，先求$c$，再用$e=\frac{c}{a}$。

		\item \textbf{场景二（定义问题）：} 题目中出现焦点、准线、距离比例时，直接使用$e=\frac{|PF|}{d(P,l)}$。

		\item \textbf{场景三（抛物线问题）：} 抛物线不要套用$e=\frac{c}{a}$，因为抛物线没有半长轴$a$和半焦距$c$这一套参数，直接记住$e=1$。

	\end{itemize}

\end{explanationbox}